\section{Arquitectura del sistema}\label{sec:arquitectura}

\subsection{Grafo de control}

La figura~\ref{fig:grafo} muestra el grafo de LangGraph.
Su estado compartido contiene la pregunta, la evidencia recuperada, el reporte del Investigador, el veredicto del Auditor, la retroalimentación, el número de rondas y la respuesta final.

\begin{figure}[ht]
\centering
\begin{tikzpicture}[
  node distance=0.9cm and 1.1cm,
  >={Stealth[length=2.2mm]},
  every node/.style={font=\small},
  agent/.style={draw, rounded corners, fill=multi!12, minimum height=1.1cm, minimum width=2.0cm, align=center},
  tool/.style={draw, fill=gray!10, minimum height=0.8cm, minimum width=1.6cm, align=center},
  term/.style={draw, ellipse, minimum height=0.7cm, align=center},
  decision/.style={draw, diamond, aspect=1.6, fill=naive!15, inner sep=1pt, align=center},
]
  \node[term] (start) {inicio};
  \node[tool, minimum width=2.0cm, above right=0.25cm and 0.55cm of start] (kb) {\code{kb}\\ ChromaDB};
  \node[tool, minimum width=2.0cm, below right=0.25cm and 0.55cm of start] (web) {\code{web}\\ DuckDuckGo};
  % El Investigador se ubica a partir del borde derecho de kb y web, así la unión no las toca.
  \coordinate (merge) at ([xshift=4mm]web.east |- start);
  \node[agent, right=4mm of merge] (inv) {Investigador\\ \footnotesize ReAct ($\leq 6$ pasos)};
  \node[decision, right=0.8cm of inv] (aud) {Auditor\\de Hechos};
  \node[agent, right=1.1cm of aud] (red) {Redactor};
  \node[term, right=0.4cm of red] (end) {fin};
  \node[tool, below=1.3cm of aud] (fb) {feedback\\ \footnotesize motivo + respuesta};
  \node[term, above=1.1cm of aud] (abs) {abstención};

  \draw[->] (start) |- (kb);
  \draw[->] (start) |- (web);
  \draw (kb.east) -| (merge);
  \draw (web.east) -| (merge);
  \draw[->] (merge) -- (inv.west);
  \draw[->] (inv) -- (aud);
  \draw[->] (aud) -- node[above, pos=0.42, font=\footnotesize] {aprobado} (red);
  \draw[->] (red) -- (end);
  \draw[->] (aud) -- node[right, font=\footnotesize] {rechazo} (fb);
  \draw[->] (fb) -| (inv);
  \draw[->] (aud) -- node[right, font=\footnotesize, align=left] {$\geq$ \code{max\_rounds}\\rechazos} (abs);
  \node[font=\footnotesize, gray, below=0.05cm of web] {en paralelo (fan-out)};
\end{tikzpicture}
\caption{Grafo cíclico de control. \code{kb} y \code{web} se consultan en paralelo antes del bucle ReAct (solo en WikiQA). Si el Auditor rechaza, el motivo y la respuesta rechazada vuelven al Investigador. Tras \code{max\_rounds}~$=3$ rechazos el sistema se abstiene.}
\label{fig:grafo}
\end{figure}

\subsection{Agentes}

\paragraph{Investigador.}
Ejecuta un bucle ReAct de hasta seis pasos.
En cada paso el LLM escribe un \emph{Thought} y una sola \emph{Action} con la forma \code{herramienta[argumento]}.
Solo el sistema escribe la \emph{Observation}: si el modelo inventa una, se descarta todo lo que viene después.
El bucle termina cuando el modelo escribe \code{Final:}.
El prompt incluye un ejemplo resuelto propio de cada tarea, la retroalimentación del Auditor y las observaciones acumuladas.

\paragraph{Auditor de Hechos.}
Es el guardrail del sistema y aplica las reglas en este orden:
\begin{enumerate}
  \item Rechaza si no hubo respuesta final dentro del límite de pasos.
  \item En código, ejecuta la respuesta junto con los \code{assert} del enunciado y la aprueba solo si pasa; esta verificación no consulta al LLM.
  \item Rechaza si la respuesta no tiene evidencia, o si incluye código que no se ejecutó tal cual.
  \item En los demás casos, le pide al LLM que responda \code{APROBADO} o \code{RECHAZADO: motivo}, con una regla específica de la tarea. Por ejemplo, en GSM8K le pide verificar que el número final responda exactamente lo que se pregunta.
\end{enumerate}

\paragraph{Redactor.}
Reformula la respuesta aprobada en una a tres oraciones, sin agregar datos ni fuentes.
En matemáticas termina con \code{Respuesta final: <número>}.
El código ya verificado se entrega sin cambios: reescribirlo con un LLM solo podría estropearlo.

\subsection{Herramientas por tarea}

Se construye un grafo por tipo de tarea con \code{build\_flow(..., task=...)}.
Cada grafo tiene sus herramientas (tabla~\ref{tab:herramientas}), su ejemplo resuelto y sus reglas de auditoría y redacción.
En GSM8K y MBPP no se ofrece la web: no aporta y podría devolver la solución publicada del propio benchmark.

\begin{table}[ht]
\centering
\caption{Herramientas disponibles para el Investigador en cada tarea.}\label{tab:herramientas}
\small
\begin{tabularx}{\linewidth}{@{}l l X@{}}
\toprule
Tarea & Herramienta & Implementación \\
\midrule
\multirow{2}{*}{WikiQA} & \code{kb[...]} & Búsqueda semántica en ChromaDB (3 fragmentos) \\
 & \code{web[...]} & DuckDuckGo Search \\
GSM8K & \code{calc[...]} & Evaluador aritmético sobre el árbol sintáctico de Python, sin \code{eval}, con límite de exponente \\
MBPP & \code{python[...]} & Proceso aislado (\code{python -I}, entorno vacío, directorio temporal) con límites de tiempo, CPU y memoria \\
\bottomrule
\end{tabularx}
\end{table}

\subsection{Concurrencia}

En WikiQA, \code{kb} y \code{web} son consultas independientes.
El grafo las lanza en paralelo con un \emph{fan-out} de LangGraph.
Las dos ramas escriben en el mismo campo del estado, que usa un reductor de concatenación, y el Investigador espera a ambas antes de empezar.
Si una rama falla, devuelve una lista vacía y no detiene a la otra.
Una prueba con una barrera de sincronización comprueba que ambas ramas se ejecutan a la vez.

\subsection{Ajustes por los límites del modelo de 7B}

Las primeras pruebas en la T4 mostraron comportamientos que se corrigieron en el diseño:
\begin{itemize}
  \item \textbf{Acciones repetidas en bucle.} Se agregaron una regla anti-repetición (repetir una \emph{Action} devuelve un aviso en vez de ejecutarla), un ejemplo resuelto y un último paso que exige \code{Final}.
  \item \textbf{Código transcrito con errores.} Si el modelo dice que su código pasó las pruebas pero no lo copia en \code{Final}, se entrega el último código que realmente ejecutó.
  \item \textbf{Relleno en la redacción.} El Redactor inventaba explicaciones que bajaban la fidelidad, así que su prompt prohíbe agregar contenido.
  \item \textbf{Idioma.} La base y la web están en inglés, y las búsquedas en español recuperaban fragmentos irrelevantes. El Investigador busca en el idioma de la pregunta.
\end{itemize}

\subsection{Arquitectura de software}

El código sigue una arquitectura hexagonal.
El núcleo (\code{domain} y \code{application}) solo depende de puertos (\code{LLMPort}, \code{KnowledgeBasePort}, \code{WebSearchPort}, \code{CodeRunnerPort}, entre otros), y los adaptadores implementan cada puerto.
Cambiar Mistral por otro LLM, o ChromaDB por FAISS, es reemplazar un adaptador.
Esto permite además probar todo el flujo sin GPU ni red con adaptadores falsos: el repositorio tiene 59 pruebas automáticas.
